% Chapter 29: The Willmore Conjecture
\let\LocalMacro\relax

\chapter{Willmore Conjecture}

\section{The Problem}

\subsection{Informal}

\textbf{Background.} Every curved surface has a natural measure of how much it bends at each point. For a donut-shaped surface (a torus), you can sum up this bending over the entire surface to get a single number called the Willmore energy. A perfectly smooth, symmetric torus has a lower energy than a wrinkled or distorted one.

\textbf{Given.} You are given any donut-shaped surface, however distorted, twisted, or self-intersecting it may be, living in three-dimensional space.

\textbf{Goal.} Determine the smallest possible Willmore energy that any donut-shaped surface can achieve, and identify which surface achieves it.

\textbf{Constraints.} The surface must be a torus (genus 1). It may self-intersect. The energy is computed from the surface's curvature at every point, averaged over the whole surface. The bound must hold for all tori, with no exceptions.

\textbf{Example.} A perfectly round, symmetric torus (like a smooth doughnut) has Willmore energy $2\pi^2 \approx 19.74$. A lumpy, distorted torus has higher energy. The question is whether $2\pi^2$ is the absolute minimum, or whether some cleverly deformed torus could achieve a lower energy.

\subsection{Formal}
\begin{conjecture}[Willmore, 1965]
For any immersed torus $T^2 \subset \mathbb{R}^3$,
\[
\mathcal{W}(T^2) \ge 2\pi^2.
\]
Equality holds for the Clifford torus in $S^3$ (stereographically projected to $\mathbb{R}^3$).
\end{conjecture}

\subsection{Historical Context}
Willmore conjectured this in 1965. It became a central problem in geometric analysis, connecting conformal geometry, integrable systems, and algebraic geometry. The conjecture resisted all attempts for 47 years.

\subsection{What Was Known}
Willmore conjectured in 1965 that the minimum energy is $2\pi^2$, achieved by the Clifford torus. Li and Yau proved in 1993 that non-embedded tori have energy at least $2\pi^2$. Ros proved the conjecture for tori of revolution in 2005. But the general case---for all tori, including self-intersecting ones---remained open for 47 years.

\subsection{Why It Matters}
The Willmore energy measures how much a surface deviates from being round---it is the natural measure of ``smoothness'' for curved surfaces. It appears in physics (the elasticity of cell membranes), computer graphics (smoothing 3D models), and biology (the shape of red blood cells). Understanding its minimum value reveals fundamental truths about the geometry of surfaces.

\subsection{Simplified Version}
For spheres (genus 0), the answer is classical: the round sphere minimizes the Willmore energy, achieving the value $4\pi$. The torus (genus 1) is the next-simplest surface, and Willmore conjectured that the Clifford torus---a perfectly symmetric torus living in 4-dimensional space---minimizes the energy among all tori. This was proved by Marques and Neves in 2014.

\section{Mathematical Theory}


\subsection{Prerequisites}
This chapter assumes familiarity with:
\begin{itemize}
\item Differential geometry of surfaces (first/second fundamental forms, mean curvature)
\item Variational calculus and geometric measure theory
\item Basic complex analysis (Riemann surfaces, uniformization)
\item Conformal geometry
\end{itemize}

\subsection{The Willmore Functional and Conformal Invariance}
The Willmore energy measures how ``bent'' a surface is. For a surface $\Sigma \subset \mathbb{R}^3$ with mean curvature $H$ and Gaussian curvature $K$:
\[
\mathcal{W}(\Sigma) = \int_\Sigma H^2 \, dA.
\]
This looks like it depends on the embedding, but a beautiful fact (Willmore, 1965) says $\mathcal{W}$ is \emph{conformally invariant}: stretching the surface conformally (like pulling rubber sheet without tearing) does not change $\mathcal{W}$. For a torus ($\chi=0$), this can be rewritten using the Gauss--Bonnet theorem ($\int_\Sigma K \, dA = 2\pi\chi(\Sigma) = 0$) as:
\[
\mathcal{W}(\Sigma) = \int_\Sigma (H^2 - K) \, dA.
\]
The integrand $H^2 - K = \frac{1}{4}(k_1 - k_2)^2 \ge 0$, where $k_1, k_2$ are the principal curvatures (the maximum and minimum bending rates at each point). This is always non-negative, and equals zero only for round spheres---so the Willmore energy is a measure of ``how far from a sphere'' the surface is.

\subsection{The Clifford Torus}
The Clifford torus in $S^3$ is
\[
T_{\text{Cl}} = \left\{ (z_1, z_2) \in \mathbb{C}^2 : |z_1| = |z_2| = \frac{1}{\sqrt{2}} \right\}.
\]
It is minimal in $S^3$. Its stereographic projection is a torus in $\mathbb{R}^3$ with Willmore energy $2\pi^2$; stereographic projection need not preserve constant mean curvature.

\section{The Solution}

\subsection{The Big Idea}

Imagine you want to prove that no donut shape can be smoother (lower Willmore energy) than the perfectly symmetric Clifford torus. The obvious approach---minimize the energy directly---fails because the space of all donut shapes is infinite-dimensional and wildly complicated. You cannot just ``optimize'' over all tori.

Marques and Neves associate a canonical family of surfaces in $S^3$ to each immersed torus. The family is built from conformal transformations of the torus, and its maximal area is bounded above by the torus's Willmore energy. They then prove that this family must contain a surface of area at least $2\pi^2$. Thus the lower bound for the family gives the desired lower bound for the original torus. The min-max argument supplies this area threshold; it does not infer a bound for every torus merely from one high-energy surface in an unrelated sweepout.

\subsection{What Was Stopping Progress}

Previous approaches tried to minimize the Willmore energy directly, but the space of immersed tori is infinite-dimensional with complicated topology. A minimizing sequence might degenerate into something that is not a torus at all. Li and Yau (1993) proved the bound for non-embedded tori---but the most important case (embedded tori) was exactly what they excluded. Ros (2005) handled tori of revolution---but those are a very special subclass. The fundamental obstacle was that variational methods could not control the topology of the minimizer: how do you guarantee a minimizing sequence stays within the right genus? Every approach was either too local (analyzing one surface at a time) or too global (trying to control the entire moduli space at once).

\subsection{The Breakthrough}

Marques and Neves used Almgren--Pitts min-max theory for the area functional in $S^3$, together with the canonical family associated to a torus. A positive-genus min-max argument gives the area threshold $2\pi^2$; comparing this threshold with the maximal area in the canonical family yields the Willmore inequality.

\subsection{How It Works}

The breakthrough of Marques and Neves (2014) \cite{marquesneves2014} was to apply \textbf{Almgren--Pitts min-max theory} for area in $S^3$ to a canonical family associated with each torus, turning a min-max area lower bound into a Willmore-energy bound.

The key ideas are:

\begin{enumerate}
\item \textbf{Min-max lower bound.} The width of the positive-genus min-max class is at least $2\pi^2$. The classification of low-index minimal surfaces in $S^3$ is a key ingredient in proving this estimate.

\item \textbf{Canonical family comparison.} For an immersed torus $\Sigma$, the canonical family consists of surfaces in $S^3$ obtained from conformal images of $\Sigma$. Its maximal area is bounded above by $\mathcal{W}(\Sigma)$, and the min-max lower bound forces that maximum to be at least $2\pi^2$.

\end{enumerate}

The proof does not directly run min-max on the fourth-order Willmore functional. It uses the area min-max theory for minimal surfaces in $S^3$, together with the conformal invariance of Willmore energy and the canonical-family construction.

\begin{theorem}[Marques-Neves \cite{marquesneves2014}]
The Willmore conjecture is true: $\mathcal{W}(T^2) \ge 2\pi^2$ for any immersed torus $T^2 \subset \mathbb{R}^3$, with equality only for the Clifford torus (up to conformal transformations).
\end{theorem}

\begin{highlightbox}
The proof is a tour de force of geometric measure theory: it combines Almgren-Pitts min-max, conformal geometry, and the classification of minimal surfaces in $S^3$ to solve a 47-year-old problem.
\end{highlightbox}

\subsection{After the Proof}

The Willmore conjecture is completely settled: Marques and Neves proved that every immersed torus in $\mathbb{R}^3$ has Willmore energy at least $2\pi^2$, with equality only for the Clifford torus. The proof is complete and no further cases or refinements of the original conjecture are needed.

The proof also strengthened the role of min-max methods in geometric analysis. In this argument, the area min-max theory is combined with conformal geometry to address a problem about Willmore energy. The connection between Willmore energy, conformal geometry, and integrable systems continues to yield new insights.

The natural open questions concern higher-genus surfaces. What is the minimum Willmore energy for an immersed surface of genus $g \ge 2$? This is known as the higher-genus Willmore problem, and the lower bound $\mathcal W\geq 2\pi^2$ holds for positive genus, while identifying minimizers for each fixed genus is a separate question. The Willmore energy of non-orientable surfaces (such as the real projective plane) is also not fully understood. Determining whether the min-max approach can be extended to these settings, or whether fundamentally new ideas are needed, is an active area of research in geometric analysis.

\section{Impact}
\begin{itemize}
\item \textbf{Geometric analysis}: Almgren-Pitts min-max became the standard tool for Willmore-type problems.
\item \textbf{Integrable systems}: Willmore surfaces are related to soliton equations; the Clifford torus corresponds to a finite-gap solution.
\item \textbf{Conformal geometry}: The Willmore energy is the natural energy for conformal immersions.
\item \textbf{Physics}: Willmore energy appears in membrane elasticity (Helfrich energy).
\end{itemize}

\section{Exercises}

\begin{exercise}[Basic]
Compute the Willmore energy of the Clifford torus in $\mathbb{R}^3$ (stereographically projected from $S^3$).
\end{exercise}

\begin{exercise}[Intermediate]
Show that the Willmore energy is conformally invariant in $\mathbb{R}^3$. Why does this fail in higher dimensions?
\end{exercise}

\begin{exercise}[Advanced]
Explain how Almgren-Pitts min-max produces a minimal surface in $S^3$ from a given sweepout of a torus.
\end{exercise}

\section{Timeline}

\begin{tabularx}{\linewidth}{lX}
\toprule
\textbf{Year} & \textbf{Event} \\ \midrule
1965 & Willmore conjectures $2\pi^2$ lower bound for tori \\ 
1993 & Li-Yau proves $\mathcal{W} \ge 2\pi^2$ for non-embedded tori \\ 
2005 & Ros proves Willmore conjecture for tori of revolution \\ 
2012--2014 & Marques--Neves prove Willmore conjecture (preprint 2012, published 2014) \\
2014 & Marques and Neves prove the Willmore conjecture \\ \bottomrule
\end{tabularx}

\section{Citations}

\begin{enumerate}
\item Marques, F. C., \& Neves, A. (2014). ``Min-max theory and the Willmore conjecture.'' Annals of Mathematics, 179(2), 683--782.
\item Willmore, T. J. (1965). ``Note on embedded surfaces.'' Anais da Academia Brasileira de Ciências, 37, 475--478.
\item Simon, L. (1993). ``Existence of surfaces minimizing the Willmore functional.'' Communications in Analysis and Geometry, 1(2), 281--326.
\end{enumerate}
